\scanpage{92}
\unsure{-- ou} seulement, \unsure{sans} doute, l'ens des pts fermés (si $S$ est de type fini sur un corps, tel $\mathbf Q$ disons). En termes de sections d'extensions, on peut dire que la section $\mathcal E_{f_K}$ de $\mathcal E_{U_K}$ sur $\mathcal E_K$ donne naissance à une famille de sections des $\mathcal E_{X_s}$ sur les $\mathcal E_{k(s)}$ pour $s\in S'$. On peut \deleted{\unsure{interpréter}} \inserted{\unsure{prolonger}} \unsure{ainsi|aussi} cette correspondance entre sections de $\mathcal E_{U_K}$ sur $\mathcal E_K$, et sections des $\mathcal E_{X_s}$ sur les $\mathcal E_{k(s)}$, en généralisant ce qu'on a vu dans le cas où $s$ est de codim $1$ : On introduit le normalisé $S'$ de $S$ dans $\overline K$, la fibre en $s$ s'identifie à l'ens des idéaux maximaux de la clôture \unsure{normale} de $\underline{\mathcal O}_{S,s}$ dans $\overline K$, et si $N_{s'}$ désigne le \deleted{\unsure{normalisateur}} \inserted{stabilisateur} dans $\mathcal E_{\overline K}$ d'un $s'$ de $S'_s$, on a encore une structure d'extension
\[
(29)\qquad
1\mathrel{\text{---}}I_{s'}\longrightarrow N_{s'}\longrightarrow \unsure{\mathcal E_{k(s)}|E_{k(s)}}\longrightarrow 1
\]
\margin{91}
